% GENERATED FILE. Edit canonical lessons, then run npm run generate:books.
% canonical-source-hash: fnv1a64:1cf08eafa31cf38a
% canonical-lessons: KA-C521-dudi, KA-C521-ustuvari, KA-C521-ottayisu, KA-C521-parisrama, KA-C521-paipoti

\chapter{Hard Work}
\label{ch:ka-hard-work}
\hlchaptermodality{521}

\begin{chapteropening}
\textbf{By the end of this chapter:} \emph{I can say to toil, oversight, to insist, hard work and rivalry.}

Complete the last lesson of chapter 521: I can say to toil, oversight, to insist, hard work and rivalry.
\end{chapteropening}

\section[\texorpdfstring{\kn{ದುಡಿ} (duḍi)}{ದುಡಿ (duḍi)}]{\kn{ದುಡಿ} (duḍi) — to toil; to labour}

\label{lesson:KA-C521-dudi}



\begin{quote}
Before the new one: say the Kannada for supervision, then the Kannada for punctuality.
\end{quote}

\subsection*{You'll want to know: \kn{ದುಡಿ}}
\textbf{\kn{ದುಡಿ}} — \emph{duḍi} — \textquotedblleft{}to toil; to labour\textquotedblright{}.

\begin{quote}
\kn{ಕುಟುಂಬಕ್ಕಾಗಿ} \kn{ಅಪ್ಪ} \kn{ಹಗಲು} \kn{ರಾತ್ರಿ} \kn{ದುಡಿದರು}.
\end{quote}

\begin{quote}
\emph{kuṭumbakkāgi appa hagalu rātri duḍidaru.}
\end{quote}

\begin{quote}
Father toiled day and night for the family.
\end{quote}

Past (I …): \kn{ದುಡಿದೆ}, \emph{duḍide}.

\kn{ದುಡಿಮೆ} is the work and what it earns.

\begin{cousinweb}
\textbf{\kn{ದುಡಿ}} (\emph{duḍi}) is native Dravidian: DEDR 3295 lists Kannada \emph{duḍi}, \textquotedblleft{}to labour, acquire by one's labour\textquotedblright{}, and \emph{duḍime}, \textquotedblleft{}gain\textquotedblright{}.
\end{cousinweb}

\begin{grammarlens}[title={What you've built}]
One more word for describing what you lived and felt.
\end{grammarlens}

\subsection*{Your turn}
\begin{itemize}
\raggedright
  \item \emph{Say it:} \emph{duḍi}
  \item \emph{Say it:} \emph{duḍi}, once more
  \item \emph{Say it:} \emph{kuṭumbakkāgi appa hagalu rātri duḍidaru.} — Father toiled day and night for the family.
  \item \emph{Recall:} say the Kannada for supervision, then the Kannada for punctuality, then say \emph{duḍi} again
\end{itemize}

\begin{tcolorbox}[breakable,colback=teal!4,colframe=teal!35!black,title={Before you move on}]
What does \kn{ದುಡಿ} mean? (\textquotedblleft{}To toil; to labour\textquotedblright{}.) Say it once more.
\end{tcolorbox}
\section[\texorpdfstring{\kn{ಉಸ್ತುವಾರಿ} (ustuvāri)}{ಉಸ್ತುವಾರಿ (ustuvāri)}]{\kn{ಉಸ್ತುವಾರಿ} (ustuvāri) — oversight; charge}

\label{lesson:KA-C521-ustuvari}



\begin{quote}
Before the new one: say the Kannada for punctuality, then the Kannada for to toil.
\end{quote}

\subsection*{You'll want to know: \kn{ಉಸ್ತುವಾರಿ}}
\textbf{\kn{ಉಸ್ತುವಾರಿ}} — \emph{ustuvāri} — \textquotedblleft{}oversight; charge\textquotedblright{}.

\begin{quote}
\kn{ಹೊಸ} \kn{ಕಟ್ಟಡದ} \kn{ಕೆಲಸದ} \kn{ಉಸ್ತುವಾರಿ} \kn{ನನ್ನ} \kn{ಅಣ್ಣನದು}.
\end{quote}

\begin{quote}
\emph{hosa kaṭṭaḍada kelasada ustuvāri nanna aṇṇanadu.}
\end{quote}

\begin{quote}
My elder brother is in charge of the new building work.
\end{quote}

\kn{ಉಸ್ತುವಾರಿ} \kn{ವಹಿಸು}, to take charge.

\begin{cousinweb}
\textbf{\kn{ಉಸ್ತುವಾರಿ}} (\emph{ustuvāri}) came into Kannada with the language of offices; its source is not settled.
\end{cousinweb}

\begin{grammarlens}[title={What you've built}]
One more word for describing what you lived and felt.
\end{grammarlens}

\subsection*{Your turn}
\begin{itemize}
\raggedright
  \item \emph{Say it:} \emph{ustuvāri}
  \item \emph{Say it:} \emph{ustuvāri}, once more
  \item \emph{Say it:} \emph{hosa kaṭṭaḍada kelasada ustuvāri nanna aṇṇanadu.} — My elder brother is in charge of the new building work.
  \item \emph{Recall:} say the Kannada for punctuality, then the Kannada for to toil, then say \emph{ustuvāri} again
\end{itemize}

\begin{tcolorbox}[breakable,colback=teal!4,colframe=teal!35!black,title={Before you move on}]
What does \kn{ಉಸ್ತುವಾರಿ} mean? (\textquotedblleft{}Oversight; charge\textquotedblright{}.) Say it once more.
\end{tcolorbox}
\section[\texorpdfstring{\kn{ಒತ್ತಾಯಿಸು} (ottāyisu)}{ಒತ್ತಾಯಿಸು (ottāyisu)}]{\kn{ಒತ್ತಾಯಿಸು} (ottāyisu) — to insist; to urge}

\label{lesson:KA-C521-ottayisu}



\begin{quote}
Before the new one: say the Kannada for to toil, then the Kannada for oversight.
\end{quote}

\subsection*{You'll want to know: \kn{ಒತ್ತಾಯಿಸು}}
\textbf{\kn{ಒತ್ತಾಯಿಸು}} — \emph{ottāyisu} — \textquotedblleft{}to insist; to urge\textquotedblright{}.

\begin{quote}
\kn{ಊಟ} \kn{ಮಾಡಿಕೊಂಡು} \kn{ಹೋಗು} \kn{ಅಂತ} \kn{ಅಜ್ಜಿ} \kn{ಒತ್ತಾಯಿಸಿದರು}.
\end{quote}

\begin{quote}
\emph{ūṭa māḍikoṇḍu hōgu anta ajji ottāyisidaru.}
\end{quote}

\begin{quote}
Grandmother insisted that I eat before going.
\end{quote}

Past (I …): \kn{ಒತ್ತಾಯಿಸಿದೆ}, \emph{ottāyiside}.

The one urged takes -\kn{ಅನ್ನು}; \kn{ಒತ್ತಾಯ} is pressure, insistence.

\begin{cousinweb}
\textbf{\kn{ಒತ್ತಾಯಿಸು}} (\emph{ottāyisu}) is usually traced to the native verb \emph{ottu}, \textquotedblleft{}to press, push\textquotedblright{} (DEDR 1021): to put pressure on someone. The middle part is not settled.
\end{cousinweb}

\begin{grammarlens}[title={What you've built}]
One more word for describing what you lived and felt.
\end{grammarlens}

\subsection*{Your turn}
\begin{itemize}
\raggedright
  \item \emph{Say it:} \emph{ottāyisu}
  \item \emph{Say it:} \emph{ottāyisu}, once more
  \item \emph{Say it:} \emph{ūṭa māḍikoṇḍu hōgu anta ajji ottāyisidaru.} — Grandmother insisted that I eat before going.
  \item \emph{Recall:} say the Kannada for to toil, then the Kannada for oversight, then say \emph{ottāyisu} again
\end{itemize}

\begin{tcolorbox}[breakable,colback=teal!4,colframe=teal!35!black,title={Before you move on}]
What does \kn{ಒತ್ತಾಯಿಸು} mean? (\textquotedblleft{}To insist; to urge\textquotedblright{}.) Say it once more.
\end{tcolorbox}
\section[\texorpdfstring{\kn{ಪರಿಶ್ರಮ} (pariśrama)}{ಪರಿಶ್ರಮ (pariśrama)}]{\kn{ಪರಿಶ್ರಮ} (pariśrama) — hard work; toil}

\label{lesson:KA-C521-parisrama}



\begin{quote}
Before the new one: say the Kannada for oversight, then the Kannada for to insist.
\end{quote}

\subsection*{You'll want to know: \kn{ಪರಿಶ್ರಮ}}
\textbf{\kn{ಪರಿಶ್ರಮ}} — \emph{pariśrama} — \textquotedblleft{}hard work; toil\textquotedblright{}.

\begin{quote}
\kn{ಅವಳ} \kn{ಯಶಸ್ಸಿಗೆ} \kn{ಅವಳ} \kn{ಪರಿಶ್ರಮವೇ} \kn{ಕಾರಣ}.
\end{quote}

\begin{quote}
\emph{avaḷa yaśassige avaḷa pariśramavē kāraṇa.}
\end{quote}

\begin{quote}
Her hard work is the reason for her success.
\end{quote}

\kn{ಪರಿ}-, all round + \kn{ಶ್ರಮ}, labour.

\begin{cousinweb}
\textbf{\kn{ಪರಿಶ್ರಮ}} (\emph{pariśrama}) is borrowed from Sanskrit \emph{pariśrama}, \textquotedblleft{}fatigue, pain, exertion\textquotedblright{} (CDIAL 7892).
\end{cousinweb}

\begin{grammarlens}[title={What you've built}]
One more word for describing what you lived and felt.
\end{grammarlens}

\subsection*{Your turn}
\begin{itemize}
\raggedright
  \item \emph{Say it:} \emph{pariśrama}
  \item \emph{Say it:} \emph{pariśrama}, once more
  \item \emph{Say it:} \emph{avaḷa yaśassige avaḷa pariśramavē kāraṇa.} — Her hard work is the reason for her success.
  \item \emph{Recall:} say the Kannada for oversight, then the Kannada for to insist, then say \emph{pariśrama} again
\end{itemize}

\begin{tcolorbox}[breakable,colback=teal!4,colframe=teal!35!black,title={Before you move on}]
What does \kn{ಪರಿಶ್ರಮ} mean? (\textquotedblleft{}Hard work; toil\textquotedblright{}.) Say it once more.
\end{tcolorbox}
\section[\texorpdfstring{\kn{ಪೈಪೋಟಿ} (paipōṭi)}{ಪೈಪೋಟಿ (paipōṭi)}]{\kn{ಪೈಪೋಟಿ} (paipōṭi) — rivalry; stiff competition}

\label{lesson:KA-C521-paipoti}



\begin{quote}
Before the new one: say the Kannada for to insist, then the Kannada for hard work.
\end{quote}

\subsection*{You'll want to know: \kn{ಪೈಪೋಟಿ}}
\textbf{\kn{ಪೈಪೋಟಿ}} — \emph{paipōṭi} — \textquotedblleft{}rivalry; stiff competition\textquotedblright{}.

\begin{quote}
\kn{ಎರಡು} \kn{ಅಂಗಡಿಗಳ} \kn{ನಡುವೆ} \kn{ತುಂಬಾ} \kn{ಪೈಪೋಟಿ} \kn{ಇದೆ}.
\end{quote}

\begin{quote}
\emph{eraḍu aṅgaḍigaḷa naḍuve tumbā paipōṭi ide.}
\end{quote}

\begin{quote}
There is fierce rivalry between the two shops.
\end{quote}

Closer and fiercer than \kn{ಸ್ಪರ್ಧೆ}.

\begin{cousinweb}
\textbf{\kn{ಪೈಪೋಟಿ}} (\emph{paipōṭi}) is built on the native \emph{pōṭi}, \textquotedblleft{}competition, rivalry\textquotedblright{}, which DEDR 4583 lists for Kannada beside Tamil \emph{pōṭṭi}; the origin of the first part, \emph{pai-}, is not settled.
\end{cousinweb}

\begin{grammarlens}[title={What you've built}]
One more word for describing what you lived and felt.
\end{grammarlens}

\subsection*{Your turn}
\begin{itemize}
\raggedright
  \item \emph{Say it:} \emph{paipōṭi}
  \item \emph{Say it:} \emph{paipōṭi}, once more
  \item \emph{Say it:} \emph{eraḍu aṅgaḍigaḷa naḍuve tumbā paipōṭi ide.} — There is fierce rivalry between the two shops.
  \item \emph{Recall:} say the Kannada for to insist, then the Kannada for hard work, then say \emph{paipōṭi} again
\end{itemize}

\begin{tcolorbox}[breakable,colback=teal!4,colframe=teal!35!black,title={Before you move on}]
What does \kn{ಪೈಪೋಟಿ} mean? (\textquotedblleft{}Rivalry; stiff competition\textquotedblright{}.) Say it once more.
\end{tcolorbox}
